% ======================================================================
\subsection{PCP-factorized self-energy: full derivation}
\label{sub:pcp_derivation}
% ======================================================================

We substitute the PCP decomposition into the NEGF self-energy and
reduce, step by step, using only universally valid identities
($e^a e^b = e^{a+b}$, linearity of sums, relabelling of dummy
indices).  No assumptions about $\vq$-mesh commensurability are
made at any point.


% -------------------------------------------------------
\subsubsection{Starting point}
% -------------------------------------------------------

\paragraph{PCP ansatz.}
The mass-weighted FC3 in compact form
($\kappa_1$ = primitive atom, $j,k$ = supercell atoms) is
\begin{equation}\label{eq:pcp_raw}
  \Psi[\kappa_1, j, k, \alpha, \beta, \gamma]
  = \sum_{\xi=1}^{N_c} \frac{\lambda_\xi}{6}
    \sum_{\sigma \in S_3} \sum_{l=0}^{N_\mathrm{cell}-1}
    m_{\sigma_1}^{\xi}[\kappa_1, l, \alpha]\;
    m_{\sigma_2}^{\xi}[j, l, \beta]\;
    m_{\sigma_3}^{\xi}[k, l, \gamma],
\end{equation}
where the mode functions are real-valued:
\begin{align}
  m_s^{\xi}[\kappa, l, \alpha]
    &= A_s^{\xi}\bigl((-l)\bmod N,\; \kappa,\; \alpha\bigr),
  \label{eq:mode_ext}\\
  m_s^{\xi}[j, l, \beta]
    &= A_s^{\xi}\bigl((t_j - l)\bmod N,\; p_j,\; \beta\bigr),
  \label{eq:mode_int}
\end{align}
with $t_j = \mathrm{cell}(j)$, $p_j = \mathrm{prim}(j)$, and
$A_s^\xi \in \RR^{N_\mathrm{cell} \times n_\mathrm{at} \times 3}$.

\paragraph{Self-energy.}
\begin{equation}\label{eq:sigma_start}
  \Sigma^{\gtrless}_{ab}(\omega;\vq)
  = \frac{i\hbar}{2N_q} \sum_{\vq'}
    \sum_{c,d,e,f}
    \Phi_{acd}(\vq',\vq{-}\vq')\;
    K_{cdfe}(\omega;\vq',\vq{-}\vq')\;
    \Phi^*_{bef}(\vq{-}\vq',\vq'),
\end{equation}
with $K_{cdfe} = \int \frac{d\omega'}{2\pi}\,
G_{cf}(\omega';\vq')\,G_{de}(\omega{-}\omega';\vq{-}\vq')$.

\paragraph{Vertex via gathering matrices.}
The Fourier-transformed vertex at $(\vq_1, \vq_2)$ is
\begin{equation}\label{eq:vertex_def}
  \Phi_{(k_1\alpha)(k_2\beta)(k_3\gamma)}(\vq_1, \vq_2)
  = \sum_{\substack{j:\, p_j = \kappa_2 \\ k:\, p_k = \kappa_3}}
    e^{-i\vq_1 \cdot \vR_{t_j}}\;
    \Psi[\kappa_1, j, k, \alpha, \beta, \gamma]\;
    e^{-i\vq_2 \cdot \vR_{t_k}},
\end{equation}
where $\vR_t$ denotes the lattice vector of cell~$t$.  This is
exact for any $\vq_1, \vq_2$---no approximations.


% -------------------------------------------------------
\subsubsection{Step 1: Substitute PCP into the vertex}
\label{sssub:step1}
% -------------------------------------------------------

Insert \cref{eq:pcp_raw} into \cref{eq:vertex_def}:
\begin{equation}\label{eq:vertex_pcp_raw}
  \Phi_{acd}(\vq_1, \vq_2)
  = \sum_\xi \frac{\lambda_\xi}{6} \sum_\sigma \sum_l
    \underbrace{m_{\sigma_1}^{\xi}[\kappa_1, l, \alpha]}_{
      \equiv\; w^{(\sigma_1,\xi,l)}_a \;\in\;\RR}
    \cdot
    \underbrace{
      \sum_{j:\,p_j=\kappa_2}
      e^{-i\vq_1\cdot\vR_{t_j}}\;
      m_{\sigma_2}^{\xi}[j, l, \beta]
    }_{\equiv\; h_{\sigma_2}^{\xi}(\vq_1, l)_c \;\in\;\CC}
    \cdot
    \underbrace{
      \sum_{k:\,p_k=\kappa_3}
      e^{-i\vq_2\cdot\vR_{t_k}}\;
      m_{\sigma_3}^{\xi}[k, l, \gamma]
    }_{\equiv\; h_{\sigma_3}^{\xi}(\vq_2, l)_d \;\in\;\CC}\,.
\end{equation}
We have defined two objects:
\begin{itemize}
  \item $w_a^{(s,\xi,l)} = A_s^\xi\bigl((-l)\bmod N,\; \kappa_a,\; \alpha_a\bigr)$:
    a \textbf{real scalar} (external mode weight).
  \item $h_s^{\xi}(\vq, l)_c$: a \textbf{complex $M$-vector} (shifted
    mode FT), defined as
    \begin{equation}\label{eq:shifted_ft}
      \boxed{
      h_s^{\xi}(\vq, l)_{(\kappa\beta)}
      = \sum_{t=0}^{N_\mathrm{cell}-1}
        e^{-i\vq\cdot\vR_t}\;
        A_s^{\xi}\bigl((t - l)\bmod N,\; \kappa,\; \beta\bigr).
      }
    \end{equation}
\end{itemize}
The vertex is therefore a sum of $6 N_c N_\mathrm{cell}$ rank-1 terms
in the $(c, d)$-indices:
\begin{equation}\label{eq:vertex_rank1}
  \Phi_{acd}(\vq_1, \vq_2)
  = \sum_\xi \frac{\lambda_\xi}{6} \sum_\sigma \sum_l
    w_a^{(\sigma_1,\xi,l)}\;
    h_{\sigma_2}^{\xi}(\vq_1, l)_c\;
    h_{\sigma_3}^{\xi}(\vq_2, l)_d\,.
\end{equation}
\emph{No approximations have been made.}  Eq.~\eqref{eq:shifted_ft}
is a standard Fourier sum, valid at any~$\vq$.


% -------------------------------------------------------
\subsubsection{Step 2: Substitute into the self-energy}
\label{sssub:step2}
% -------------------------------------------------------

Insert \cref{eq:vertex_rank1} (and its conjugate) into
\cref{eq:sigma_start}.  The left vertex contributes
\begin{equation}
  \Phi_{acd} = \sum_{\xi,\sigma,l}
    \frac{\lambda_\xi}{6}\;
    w_a^{(\sigma_1,\xi,l)}\;
    h_{\sigma_2}^{\xi}(\vq', l)_c\;
    h_{\sigma_3}^{\xi}(\vq{-}\vq', l)_d\,,
\end{equation}
and the right vertex (with $\vq_1 \leftrightarrow \vq_2$ swapped,
then complex-conjugated):
\begin{equation}
  \Phi^*_{bef} = \sum_{\xi',\sigma',l'}
    \frac{\lambda_{\xi'}}{6}\;
    w_b^{(\sigma'_1,\xi',l')}\;
    h_{\sigma'_2}^{\xi'\,*}(\vq{-}\vq', l')_e\;
    h_{\sigma'_3}^{\xi'\,*}(\vq', l')_f\,.
\end{equation}

The contraction over internal indices $(c,d,e,f)$ with $K_{cdfe}$
involves two Green's functions at different momenta:
\begin{align}\label{eq:contraction_step}
  &\sum_{c,d,e,f}
    h_{\sigma_2}(\vq', l)_c\;
    h_{\sigma_3}(\vq{-}\vq', l)_d\;
    G_{cf}(\omega';\vq')\;
    G_{de}(\omega{-}\omega';\vq{-}\vq')\;
    h^*_{\sigma'_2}(\vq{-}\vq', l')_e\;
    h^*_{\sigma'_3}(\vq', l')_f
  \notag\\[6pt]
  &=
    \underbrace{
      \Bigl[\sum_{c,f}
      h_{\sigma_2}(\vq', l)_c\;
      G_{cf}(\omega';\vq')\;
      h^*_{\sigma'_3}(\vq', l')_f
      \Bigr]
    }_{\displaystyle
      g_{\sigma_2,\sigma'_3}^{\xi\xi'}(\omega';\vq',l,l')
    }
    \;\times\;
    \underbrace{
      \Bigl[\sum_{d,e}
      h_{\sigma_3}(\vq{-}\vq', l)_d\;
      G_{de}(\omega{-}\omega';\vq{-}\vq')\;
      h^*_{\sigma'_2}(\vq{-}\vq', l')_e
      \Bigr]
    }_{\displaystyle
      g_{\sigma_3,\sigma'_2}^{\xi\xi'}(\omega{-}\omega';\vq{-}\vq',l,l')
    }\,.
\end{align}
This factorization is \textbf{exact}: $(c,f)$ both pair with
$G(\vq')$ and $(d,e)$ both pair with $G(\vq{-}\vq')$.  The two
factors are independent.

\paragraph{Scalar Green's function projections.}
Define:
\begin{equation}\label{eq:scalar_proj}
  \boxed{
  g_{ss'}^{\xi\xi'}(\omega;\vq,l,l')
  = h_s^{\xi}(\vq, l)^\dagger\;
    G^{\gtrless}(\omega;\vq)\;
    h_{s'}^{\xi'}(\vq, l').
  }
\end{equation}
These are \textbf{scalar} functions of $(\omega, \vq)$, parameterized
by $(s,s',\xi,\xi',l,l')$.


% -------------------------------------------------------
\subsubsection{Step 3: Full PCP self-energy}
\label{sssub:step3}
% -------------------------------------------------------

Combining Steps 1--2:
\begin{equation}\label{eq:sigma_pcp_full}
  \boxed{
  \Sigma^{\gtrless}_{ab}(\omega;\vq)
  = \frac{i\hbar}{2N_q}
    \sum_{\vq'}
    \sum_{\xi,\xi'} \frac{\lambda_\xi\lambda_{\xi'}}{36}
    \sum_{\sigma,\sigma' \in S_3}
    \sum_{l,l'=0}^{N_\mathrm{cell}-1}
    w_a^{(\sigma_1,\xi,l)}\; w_b^{(\sigma'_1,\xi',l')}
    \;\mathrm{IFFT}_\omega\!\Bigl\{
      g_{\sigma_2,\sigma'_3}^{\xi\xi'}(\vq',l,l')\;
      g_{\sigma_3,\sigma'_2}^{\xi\xi'}(\vq{-}\vq',l,l')
    \Bigr\}.
  }
\end{equation}
This is exact at arbitrary $\vq$, with no commensurability
requirement.  Compared to the ``standard'' PCP expression
(which uses mode FTs $\hat f_s^\xi(\vq)$ without the $l$-index),
the shifted-FT version carries an additional double sum over
$(l, l')$.


% -------------------------------------------------------
\subsubsection{Step 4: Optimal evaluation order and cost}
\label{sssub:step4}
% -------------------------------------------------------

We now determine the optimal summation order to minimize
floating-point operations and memory.  Let:
\begin{equation*}
  M = 3 n_\mathrm{at},\quad
  N_q = n_{q_x} n_{q_y},\quad
  N_\tau = 2(n_\mathrm{low} + N_\omega),\quad
  L = N_\mathrm{cell}.
\end{equation*}

\paragraph{Phase 0: Precompute the shifted-FT modes $h$.}

For each $\vq$-point, compute
$h_s^{\xi}(\vq, l) \in \CC^M$ for all $(s, \xi, l)$.

Naively, \cref{eq:shifted_ft} requires $L$ multiplications and
additions per DOF per $(s, \xi, l, \vq)$: total
$3 N_c L^2 M N_q$.

\emph{Faster via circulant structure.}
For fixed $\vq$, define the circulant matrix
$\Phi_{lt} = e^{-i\vq\cdot\vR_{(t+l)\bmod N}}$.  Then
$h_s^{\xi}(\vq, l)_{(\kappa\alpha)}
= \sum_t \Phi_{lt}\, A_s^\xi(t, \kappa, \alpha)$.
Since $\Phi$ is circulant (its $(l,t)$-entry depends only on
$(t + l) \bmod N$), the product $\Phi \cdot A$ is a circular
convolution, computable via FFT in $\cO(L \log L)$ per column.
Total:
\begin{equation}\label{eq:cost_h}
  \cW_h = 3 N_c \cdot M \cdot L\log L \cdot N_q.
\end{equation}
\emph{Storage}: $h$ has $L \times N_q \times 3 \times N_c \times M$
complex entries.  For reuse across SCBA iterations, store
$h$ once ($\vq$-dependent but $\omega$-independent).

\paragraph{Phase 1: Scalar projections $g$.}

For each $(\vq, \omega)$: given $G(\omega;\vq) \in \CC^{M \times M}$
and $h_s^\xi(\vq, \cdot) \in \CC^{L \times M}$ (for each $s, \xi$),
compute:
\begin{equation}\label{eq:proj_matmul}
  g_{ss'}^{\xi\xi'}(\omega;\vq,l,l')
  = \bigl[h_s^{\xi}(\vq)^\dagger\; G(\omega;\vq)\;
    h_{s'}^{\xi'}(\vq)\bigr]_{ll'}
  \;\in\; \CC^{L \times L}.
\end{equation}
This is a matrix triple product $(L \times M)(M \times M)(M \times L)$.
\begin{itemize}
  \item Compute $G \cdot h_{s'}^{\xi'} \in \CC^{M \times L}$:
    cost $M^2 L$ per $(s', \xi')$.
  \item Left-multiply by $h_s^{\xi\,\dagger} \in \CC^{L \times M}$:
    cost $L^2 M$ per $(s, \xi, s', \xi')$.
\end{itemize}

We can \textbf{reuse} $G h$ across left-multipliers.  For each
$(\vq, \omega)$:
\begin{enumerate}
  \item Form $\tilde{G}_{s'\xi'} = G\; h_{s'}^{\xi'}
    \in \CC^{M \times L}$ for all $3 N_c$ combinations.
    Cost: $3 N_c \cdot M^2 L$.
  \item For each $(s, \xi)$: $g_{ss'}^{\xi\xi'} =
    h_s^{\xi\dagger}\, \tilde{G}_{s'\xi'} \in \CC^{L\times L}$.
    Cost: $9 N_c^2 \cdot L^2 M$.
\end{enumerate}
Total per $(\vq, \omega)$:
\begin{equation}\label{eq:cost_proj}
  \cW_\mathrm{proj}^{(\vq,\omega)}
  = 3 N_c M^2 L + 9 N_c^2 L^2 M.
\end{equation}
Over all $(\vq, \omega)$: multiply by $N_q \cdot N_\tau$.

\emph{Storage}: $g$ has $9 N_c^2 \cdot L^2$ scalars per
$(\vq, \omega)$.  We only need $g$ at the current $\omega$
for the FFT step, so $g$ can be streamed.

\paragraph{Phase 2: FFT convolution.}

For each $(\vq, \vq')$ pair: form the product
$\hat{g}_1(\vq') \cdot \hat{g}_2(\vq{-}\vq')$ (in the $\omega$
FFT domain) and IFFT.  The permutation sum over $S_3 \times S_3$
groups 36 products per $(\xi,\xi',l,l')$ into 9 output bins
$(s_1, s_1')$.

Concretely, for each $(\xi, \xi', l, l')$:
\begin{equation}
  C^{(\xi\xi')}_{s_1 s_1'}(\omega; \vq, l, l')
  = \mathrm{IFFT}_\omega\!\Biggl\{
    \sum_{\substack{(s_1,s_2,s_3)\in S_3\\(s_1',s_2',s_3')\in S_3}}
    \hat g_{s_2,s_3'}^{\xi\xi'}(\vq',l,l')\;
    \hat g_{s_3,s_2'}^{\xi\xi'}(\vq{-}\vq',l,l')
  \Biggr\}.
\end{equation}
Cost per $(\vq, \vq', \xi, \xi', l, l')$:
36 complex multiplications of length-$N_\tau$ arrays $+$
9 IFFTs of length~$N_\tau$.
Total:
\begin{equation}\label{eq:cost_conv}
  \cW_\mathrm{conv}
  = N_q^2 \cdot N_c^2 \cdot L^2 \cdot
    \bigl(36\, N_\tau + 9\, N_\tau \log N_\tau\bigr).
\end{equation}

\paragraph{Phase 3: Accumulation (outer product).}

For each $(\vq, \vq', \xi, \xi', l, l')$ we have
$C_{s_1 s_1'}(\omega) \in \CC$ at $N_\omega$ frequencies.
Accumulate into $\Sigma(\omega;\vq) \in \CC^{M \times M}$:
\begin{equation}\label{eq:accumulate}
  \Sigma_{ab}(\omega;\vq)
  \;\mathrel{+}=\;
  \frac{\lambda_\xi\lambda_{\xi'}}{36}\;
  w_a^{(s_1,\xi,l)}\;
  C_{s_1 s_1'}(\omega;\vq,l,l')\;
  w_b^{(s_1',\xi',l')}\,.
\end{equation}
For fixed $(\xi, \xi', l, l', s_1, s_1')$, the outer product
$w_a \cdot w_b$ can be precomputed as an $M \times M$ matrix and
reused across $\omega$ and $\vq'$:
\begin{equation}
  W^{(\xi\xi')}_{s_1 s_1', l l'}
  = \frac{\lambda_\xi\lambda_{\xi'}}{36}\;
    w^{(s_1,\xi,l)} \otimes w^{(s_1',\xi',l')}
  \;\in\; \RR^{M \times M}.
\end{equation}
Then \cref{eq:accumulate} is an $M \times M$ rank-1 update weighted
by a scalar.  The sum over $(l, l')$ can be done \emph{before}
the $\vq'$-sum by accumulating $C$ over $\vq'$ first.

\emph{Optimal order}: for each $(\vq, \xi, \xi', s_1, s_1')$:
\begin{enumerate}
  \item Sum $C$ over $\vq'$: gives
    $\bar{C}_{s_1 s_1'}^{(\xi\xi')}(\omega;\vq,l,l')
    = \sum_{\vq'} C(\omega;\vq,l,l')$.
    Size: $L^2 \times N_\omega$.
  \item Contract with $W$:
    $\Sigma_{ab} \mathrel{+}= \sum_{l,l'} W_{ll'}^{ab}\;
    \bar{C}(\omega,l,l')$.
\end{enumerate}
Cost of step~2 per $(\vq, \xi, \xi', s_1, s_1', \omega)$:
$L^2$ multiply-adds for the scalar sum, then one $M^2$-cost
rank-1 update.  Total for the accumulation:
\begin{equation}\label{eq:cost_accum}
  \cW_\mathrm{accum}
  = 9 N_c^2 \cdot N_\omega \cdot N_q \cdot (L^2 + M^2).
\end{equation}
(The $L^2$ term contracts $(l,l')$; the $M^2$ term is the outer
product.)


% -------------------------------------------------------
\subsubsection{Step 5: Total cost and comparison}
\label{sssub:step5}
% -------------------------------------------------------

\paragraph{PCP shifted-FT cost per SCBA iteration.}
\begin{equation}\label{eq:cost_pcp_total}
  \cW_\mathrm{PCP}
  = \underbrace{N_q N_\tau (3 N_c M^2 L + 9 N_c^2 L^2 M)}_{\text{projections}}
  + \underbrace{N_q^2 N_c^2 L^2 \cdot 9 N_\tau \log N_\tau}_{\text{FFT conv.}}
  + \underbrace{9 N_c^2 N_\omega N_q (L^2 + M^2)}_{\text{accumulation}}.
\end{equation}
Phase~0 (computing $h$) is done once and reused across iterations.

\paragraph{Dense cost.}
\begin{equation}\label{eq:cost_dense}
  \cW_\mathrm{dense}
  = N_q^2 (M^4 N_\tau \log N_\tau + M^5 N_\omega).
\end{equation}

\paragraph{Separable (SVD) cost.}
\begin{equation}\label{eq:cost_sep}
  \cW_\mathrm{sep}
  = N_q N_\tau (R M^2 + R^2 M) + N_q^2 R^2 N_\tau \log N_\tau
  + R^2 N_\omega N_q M^2.
\end{equation}

\paragraph{Comparison for Si ($M = 6$, $L = 8$, $N_q = 16$,
  $N_\tau = 40$, $N_\omega = 101$).}

\begin{center}
\renewcommand{\arraystretch}{1.3}
\begin{tabular}{l@{\qquad}rrrr}
  \toprule
  & Dense & Sep.\ $R{=}24$ & PCP $N_c{=}24$ & PCP $N_c{=}4$ \\
  \midrule
  Projections
    & --- & $3.7 \times 10^6$ & $5.3 \times 10^7$ & $4.1 \times 10^6$ \\
  FFT conv.
    & $4.4 \times 10^7$ & $3.8 \times 10^7$ & $5.5 \times 10^8$ & $1.5 \times 10^7$ \\
  Accumulation
    & $4.4 \times 10^6$ & $8.4 \times 10^6$ & $9.5 \times 10^7$ & $3.5 \times 10^6$ \\
  \midrule
  \textbf{Total}
    & $4.9 \times 10^7$ & $4.9 \times 10^7$ & $7.0 \times 10^8$ & $2.3 \times 10^7$ \\
  \textbf{Ratio to dense}
    & $1.0\times$ & $1.0\times$ & $14\times$ & $\mathbf{0.47\times}$ \\
  \bottomrule
\end{tabular}
\end{center}

At full rank ($N_c = 24$), the $L^2 = 64$ overhead makes PCP
$14\times$ slower.  But the whole point of PCP is \emph{compression}:
at $N_c = 4$ (FC3 error $\sim 3\%$), PCP is $\mathbf{2\times}$
faster than dense and comparable to the SVD.


% -------------------------------------------------------
\subsubsection{Step 6: Memory layout and reuse}
\label{sssub:memory}
% -------------------------------------------------------

\paragraph{Reuse across SCBA iterations.}
The shifted-FT modes $h_s^\xi(\vq, l)$ and external weights
$w_a^{(s,\xi,l)}$ depend only on the PCP fit and the $\vq$-mesh.
They are computed \emph{once} before the SCBA loop and reused in
every iteration.  Cost: $\cW_h$ (negligible).

The outer-product matrices $W^{(\xi\xi')}_{s_1 s_1', ll'}$ are
also iteration-independent.

\paragraph{Memory for $h$.}
$L \times N_q \times 3 \times N_c \times M$ complex doubles.
For Si: $8 \times 16 \times 3 \times 24 \times 6 \times 16\,\text{B}
= 5.3\,\text{MB}$.  At $N_c = 4$: $0.9\,\text{MB}$.

\paragraph{Memory for $g$.}
$9 N_c^2 \times L^2$ scalars per $(\vq, \omega)$.  At $N_c = 24$:
$9 \times 576 \times 64 = 331{,}776$ complex doubles $= 5.1\,\text{MB}$.
At $N_c = 4$: $9{,}216$ per $(\vq,\omega)$---trivial.

Since $g$ is consumed immediately in the FFT step, it needs
storage for only one $\vq$-point at a time (streamed).

\paragraph{Memory for $W$.}
$9 N_c^2 \times L^2 \times M^2$ reals.
At $N_c = 4$: $9 \times 16 \times 64 \times 36 = 331{,}776$
doubles $= 2.5\,\text{MB}$---feasible to precompute.
At $N_c = 24$: $9 \times 576 \times 64 \times 36 = 11.9\times10^6$
doubles $= 95\,\text{MB}$---large but feasible; alternatively,
compute the $(l,l')$-contracted accumulation on-the-fly.


% -------------------------------------------------------
\subsubsection{Step 7: Relation to the ``standard'' PCP kernel}
\label{sssub:relation}
% -------------------------------------------------------

The ``standard'' PCP kernel uses the simple mode FT
\begin{equation}
  \hat{f}_s^\xi(\vq)_a
  = \sum_t e^{-i\vq \cdot \vR_t}\, A_s^\xi(t, \kappa, \alpha),
\end{equation}
which is the special case $l = 0$ of the shifted FT:
$\hat{f}_s^\xi(\vq) = h_s^\xi(\vq, l{=}0)$.

When the $l$-sum can be collapsed into a phase factor, i.e., when
\begin{equation}\label{eq:shift_identity}
  h_s^\xi(\vq, l) = e^{-i\vq\cdot\vR_l}\; \hat{f}_s^\xi(\vq)
  \qquad\text{for all } l,
\end{equation}
the sums over $l$ and $l'$ in \cref{eq:sigma_pcp_full} collapse:
the phases combine with $w$ into the standard FT $\hat{f}$, and
$g$ loses its $(l,l')$-dependence.  This recovers the standard
$N_c^2$-cost kernel without the $L^2$ overhead.

Eq.~\eqref{eq:shift_identity} holds when $\vq \cdot N_i \va_i
\in 2\pi\ZZ$ for all lattice directions.  The derivation in
\cref{sssub:step1}--\cref{sssub:step3} does \emph{not} rely on
this identity and is therefore valid at arbitrary~$\vq$.


% -------------------------------------------------------
\subsubsection{Summary}
\label{sssub:summary}
% -------------------------------------------------------

\begin{enumerate}
  \item The PCP self-energy
    \cref{eq:sigma_pcp_full} is \textbf{exact at any $\vq$} when
    formulated with shifted-FT mode vectors $h_s^\xi(\vq, l)$.

  \item The cost overhead versus the standard PCP kernel is a
    factor of $L^2 = N_\mathrm{cell}^2$ in the scalar projections
    and FFT convolution steps.

  \item At low PCP rank ($N_c \ll M^2/3$), even with the $L^2$
    overhead, the factorized kernel is faster than the dense
    kernel.  For Si at $N_c = 4$: $\sim 2\times$ speedup.

  \item All $\vq$-independent quantities ($h$, $w$, $W$) are
    precomputed once and reused across SCBA iterations, giving
    further amortization.
\end{enumerate}
